%======================================================================
% Section 13: Applications of Multiphysics Wave Mechanics
%======================================================================
\section{Applications of Multiphysics Wave Mechanics}
\label{sec:applications}

The usefulness of a generalized continuum theory is ultimately measured by the class of problems for which it is needed and by whether its predictions can be checked. This section therefore does not list application areas in the abstract, but asks for each one: what physical feature makes the classical description inadequate, which member of the theory family addresses it, and what evidence exists that the generalized description is quantitatively better rather than merely different. The classification of evidence levels used in Table~\ref{tab:evidence_map} follows that logic: \emph{High} means that the effect is experimentally established and that generalized theory is required to reproduce it; \emph{Medium} means that the effect is established but generalized theory provides an improved rather than essential description; \emph{Low} means that the application is presently served by classical theory and generalized theory is speculative.

\begin{figure}[htbp]
\centering
\includegraphics[width=0.95\textwidth]{fig16-observation-workflow.pdf}
\caption{From a generalized continuum model to a testable prediction. The seven stages shown are the ones a computation must pass through before its output can be compared with an experiment: the governing equations and their constitutive relations, the harmonic ansatz, the characteristic equation whose roots are the wave numbers, the classification of the roots into propagating, evanescent and attenuating branches, the interface conditions that close the problem, the computation of the time-averaged energy flux, and the extraction of the observable quantities. The three verification checkpoints on the right --- the provenance of every equation, the statement of every assumption, and the specification of a validity domain --- are the point at which most of the reviewed literature is weakest, and they are the reason the evidence levels of Table~\ref{tab:evidence_map} are not uniformly high.}
\label{fig:observation_workflow}
\end{figure}

\subsection{Smart structures and adaptive materials}
\label{sec:13.1}

Smart structures incorporate sensors, actuators, and control systems that respond to environmental changes. Piezothermoelastic materials are used in:
\begin{itemize}
\item Adaptive aerospace structures (wing shaping, vibration control),
\item Active vibration suppression in precision machinery,
\item Structural health monitoring with embedded piezoelectric sensors.
\end{itemize}

\subsubsection{Evidence}
The coupling that makes these systems work is the same coupling that the generalized theories formalize: the piezoelectric tensor generates an electric field from mechanical strain and a strain from an electric field, and the thermal expansion coefficient couples temperature change to stress. In a layered device under a moving heat source, the three fields are not independent, and the work of Abouelregal, Askar, Marin and Mohamed (2023) \cite{Abouelregal2023} is representative of the current treatment: a functionally graded rotating rod is taken to be electrically unconnected at its ends and thermally isolated, and the coupled field response is computed in the Laplace domain for a modified Lord--Shulman model with a memory-dependent derivative. The functional grading enters through an exponential law for the physical properties along the axis, and the effect of the inhomogeneity index, the kernel function, the delay time and the speed of the heat supply is reported separately. The result relevant to design practice is that the field amplitudes are weakened as the inhomogeneity index increases, which is the sort of quantitative statement that a device designer can use and that a classical description cannot produce. \emph{Evidence level: Medium} --- the mechanism is experimentally established in piezoelectric composites, but the specific memory-dependent formulation is currently a theoretical description of a measured effect.

\subsection{Nondestructive evaluation and sensing}
\label{sec:13.2}

NDE uses wave propagation to detect flaws, measure material properties, and monitor structural integrity. Applications include:
\begin{itemize}
\item Ultrasonic testing for crack detection,
\item Thermal wave imaging for subsurface defect detection,
\item Acoustic emission monitoring for damage assessment.
\end{itemize}

\subsubsection{Evidence}
Thermal-wave inspection is the clearest case in which a generalized theory is not a refinement but a requirement. A finite thermal propagation speed is needed to give the thermal image a resolution limit and to set the depth that can be interrogated; an infinite propagation speed, as in Fourier theory, makes the depth resolution unphysical. The classical wave-propagation framework for elastic solids, including the treatment of anisotropy and the interaction of incident, reflected and surface waves, is set out by Achenbach (1973) \cite{Achenbach1973} and remains the reference for the elastic part of the problem. The thermal part is where the generalizations enter, and the balance of evidence assembled for second sound by Joseph and Preziosi \cite{Joseph1989,Joseph1990} is what distinguishes the finite-speed descriptions from the classical one. \emph{Evidence level: High} for ultrasonic NDE, where classical elasticity suffices; \emph{Medium} for thermal-wave imaging, where generalized theory is required in principle but the quantitative inversion of defect depth in practice still uses simplified models.

\subsection{Energy harvesting and conversion}
\label{sec:13.3}

Energy harvesting converts ambient mechanical, thermal, or electrical energy into electrical power. Piezothermoelectric energy harvesters use:
\begin{itemize}
\item Vibration-induced piezoelectric power generation,
\item Thermal gradient-driven thermoelectric conversion,
\item Coupled thermo-piezoelectric conversion.
\end{itemize}

\subsubsection{Evidence}
The theoretical difficulty in this area is specific and worth stating: the thermopower delivered by a piezoelectric in a temperature gradient competes with the direct piezoelectric power generated by the same mechanical deformation, and the competition depends on the thermal relaxation time of the material. Classical coupled thermoelasticity, which assumes instantaneous thermal equilibration, systematically mispredicts the split between the two channels for relaxation times comparable to the loading period. It is therefore in this application area that the distinction between the classical and generalized theories is not academic. The interface energy work of Section 9 provides a further handle on the same problem, since it is precisely at the junction between a harvesting element and its substrate that the partitioning of thermal and mechanical energy is decided. \emph{Evidence level: High} for piezoelectric harvesting in general; \emph{Medium} for the coupled thermo-piezoelectric channel, for which the theory is mature and the experimental separation of the channels is difficult.

\subsection{Geophysics and seismology}
\label{sec:13.4}

Geophysical applications include:
\begin{itemize}
\item Seismic wave propagation in the Earth's crust,
\item Thermal wave propagation in geothermal systems,
\item Poroelastic wave propagation in saturated reservoirs.
\end{itemize}

\subsubsection{Evidence}
Seismic wave theory is well established \cite{Achenbach1973}; poroelastic wave propagation has strong theoretical and experimental support. Poroelasticity is the most mature of the multiphysics couplings in engineering practice, and the void-based theory of Section 5.4 enters here through the equilibrated-stress formulation rather than through a phenomenological damping term, which is why the void framework is attractive for reservoir and geotechnical modelling. \emph{Evidence level: High.}

\subsection{Biomedical applications (ultrasound, thermal therapy)}
\label{sec:13.5}

Biomedical applications include:
\begin{itemize}
\item Medical ultrasound imaging and therapy,
\item Thermal ablation for cancer treatment,
\item Photo\-acoustic imaging,
\item Temperature-sensitive drug delivery.
\end{itemize}

\subsubsection{Evidence}
Medical ultrasound is established; evidence for thermal ablation is growing, while photo\-acoustic imaging remains an active research area. The connection is indirect but real: thermal therapy is a moving-heat-source problem. Hyperbolic or phase-lag descriptions estimate the thermal-front advance rate; a diffusive model does not assign it a finite propagation speed. \emph{Evidence level: High} for ultrasound; \emph{Medium} for thermal ablation and photo\-acoustic imaging.

\subsection{Nanoscale and semiconductor devices}
\label{sec:13.6}

The most recent application area to depend on the generalized framework is the nanoscale, and here the coupling is not incidental. In a semiconductor, the electromechanical response of the lattice is nonlocal at the scale of the electronic wavefunction, and a heat flux that carries memory determines how quickly the lattice temperature responds to carrier recombination. Gupta, Hussain, Yaylac{\i s} and Das (2026) \cite{Gupta2026HeatMap} combine a nonlocal Kelvin--Voigt piezoelectric semiconductor with a memory-dependent heat flux and map the resulting wave dynamics, which illustrates both the reach and the difficulty of the class: the model is well posed and produces interpretable fields, but each added mechanism contributes parameters that must be fixed independently.

\subsubsection{Evidence}
Nonlocal elasticity is not optional at these scales. Maranganti and Sharma (2007) \cite{Maranganti2007} showed that for common materials the length below which classical elasticity ceases to be accurate is of the order of tens of nanometres for polythene and single-digit nanometres for silica, and Maranganti, Sharma and Sharma (2006) \cite{Maranganti2006} showed that electromechanical coupling in non-piezoelectric dielectrics is a nanoscale effect. Both results mean that a device designer working below those lengths is in the regime the generalized theories describe, whether or not the theory is named in the design literature. \emph{Evidence level: Medium} --- the size dependence is experimentally established at the level of stiffness and of electromechanical response, and generalized theory is required in principle, but full multiphysics nonlocal thermoelastic models are at the level of validated formalisms rather than deployed design tools.

\subsection{Applications the generalized framework does not yet serve}
\label{sec:13.7}

An evidence map is only as useful as its negative entries, and the applications in which the generalized framework is currently \emph{not} required deserve the same treatment as those in which it is. Four are worth recording.

\textbf{Ordinary engineering thermoelasticity.} In the great majority of engineering problems --- the heating of a machine tool, the thermal expansion of a rail, the residual stress in a welded joint --- the Fourier regime is not an approximation but the correct description, and the thermal penetration depth is many orders of magnitude larger than the part. The generalized theories add parameters without adding accuracy. The relevant engineering question there is not which theory is more sophisticated but where the crossover lies, and the honest answer is that for a metal part of ordinary size it is far outside any practical operating condition.

\textbf{Long-wavelength geophysics.} Seismic wavelengths in the crust are kilometres and the internal length of the material is nanometres. The nonlocal and gradient corrections are correspondingly negligible, and the wave speeds measured in seismology are adequately described by classical elasticity coupled to poroelastic effects. Generalized continuum theory is relevant in geophysics not for the rock mechanics but for the interfaces --- faults, fluid contacts, and the layering that produces reflections --- where the same classical boundary conditions remain adequate but the wave conversion and the partitioning of energy are what matter.

\textbf{Long-duration thermal processes.} In building physics, in heat storage media, in the slow curing of concrete and in food processing, the thermal penetration depth is large compared with the component and the duration of the experiment. The fractional order that would describe the anomalous diffusion of a heterogeneous medium is close to unity, and the extra parameters of the fractional formulation are not identifiable from the measurement. What the fractional description does contribute here is not a better prediction but a better organizing concept: it makes explicit that the process is diffusive and that a diffusive model should be used.

\textbf{Coupled motion in structural composites.} Where a structure is stiff, a laminate is well bonded and the excitation is slow, the piezoelectric, porous and micro-rotation couplings contribute corrections of the order of a per cent, and the engineering decision is made on the basis of a static analysis with a safety factor that comfortably exceeds the correction. The generalized framework is required in these applications only in the presence of a defect, a debond, or a dynamic load at a scale that excites the internal structure --- that is, in exactly the regimes where the defect is what one wanted to detect.

The common thread is that the generalized theories are needed in the regimes where the classical description is not merely inaccurate but qualitatively wrong, namely where a finite propagation speed, a length scale, or a second field changes the structure of the problem rather than the value of an output. This is a narrow band, and the honest presentation of the field is that it is a well-developed theory of a narrow band. Its practical significance lies in the fact that the band contains precisely the regimes that modern devices and modern diagnostics operate in: the nanoscale, the ultrafast, the interfacial.


\begin{table}[htbp]
\centering
\small
\caption{Application areas discussed in the inherited review and evidence questions for future checking. This table assigns no validation-strength score.}
\label{tab:T8-applications}
\begin{tabularx}{\textwidth}{@{}lXX@{}}
\toprule
\textbf{Application area} & \textbf{Model families mentioned} & \textbf{Evidence audit required} \\
\midrule
Non-destructive evaluation & Elastic, thermoelastic, and interface-wave models & Identify experiments, material data, uncertainty, and independent replication \\
Geophysics & Thermoelastic and porous-media models & Separate field observations from model-generated examples; verify parameter provenance \\
MEMS/NEMS & Nonlocal and size-dependent models & Check specimen-scale data and boundary conditions against the cited studies \\
Biomedical / thermal therapy & Bioheat and coupled thermoelastic models & Verify clinical versus in-vitro/in-silico evidence and applicable safety/measurement claims \\
Energy harvesting & Piezoelectric and thermal-gradient models & Distinguish analytical efficiency from measured prototype performance \\
Smart structures & Active and coupled-field models & Check experimental setup, calibration, uncertainty, and comparison baseline \\
\bottomrule
\multicolumn{3}{@{}l}{\footnotesize Literature examples are source leads only; no application-level validation score was independently assigned.}\\
\end{tabularx}
\end{table}


\begin{figure}[htbp]
\centering
\includegraphics[width=0.95\textwidth]{fig17-application-evidence.pdf}
\caption{Application evidence map: peer-reviewed evidence levels across applications and theory families (Section 13.8).}
\label{fig:evidence_map}
\end{figure}

\subsection{Evidence map: which applications have peer-reviewed support}
\label{sec:13.8}

\begin{table}[htbp]
\centering
\small
\caption{Evidence map of applications across theory families, with the mechanism that makes a generalized description necessary and the supporting literature.}
\label{tab:evidence_map}
\begin{tabularx}{\textwidth}{@{}p{2.2cm}X p{2.2cm}X@{}}
\toprule
\textbf{Application} & \textbf{Why classical theory is inadequate} & \textbf{Evidence level} & \textbf{Key references} \\
\midrule
Smart structures & Thermal relaxation competes with the electromechanical coupling & Medium & \cite{Abouelregal2023} \\
NDE (ultrasonic) & Classical elasticity suffices & High & \cite{Achenbach1973} \\
NDE (thermal wave) & Finite thermal wave speed needed to set depth resolution & Medium & \cite{Joseph1989,Joseph1990} \\
Energy harvesting & Relaxation time controls the split between thermal and mechanical channels & High & Section 13.3 analysis \\
Geophysics/\allowbreak seismology & Poroelastic coupling requires a genuine two-field formulation & High & \cite{Cowin1983,Iesan1986} \\
Biomedical ultrasound & Classical elasticity suffices & High & \cite{Achenbach1973} \\
Thermal ablation & Thermal front speed governs depth of deposition & Medium & Section 13.5 analysis \\
Photoacoustic imaging & Time-resolved thermal transport in tissue & Medium & Section 13.5 analysis \\
Nanoscale devices & Nonlocal length comparable to device dimension & Medium & \cite{Maranganti2007,Maranganti2006} \\
Semi\-conductors & Coupled nonlocal elasticity and memory heat flux & Medium & \cite{Gupta2026HeatMap} \\
\bottomrule
\end{tabularx}
\end{table}

\subsection{Notation summary for Section 13}

\begin{table}[htbp]
\centering
\small
\caption{Notation introduced in Section 13.}
\label{tab:notation-sec13}
\begin{tabularx}{\textwidth}{@{}p{1.8cm}X p{2.2cm}p{2.1cm}@{}}
\toprule
\textbf{Symbol} & \textbf{Meaning} & \textbf{Units} & \textbf{Equation} \\
\midrule
$\chi$ & Inhomogeneity index of the functionally graded material & -- & Section 13.1 \\
$\tau$ & Thermal relaxation time & s & Section 13.3 \\
$v$ & Heat supply speed & m s$^{-1}$ & Section 13.1 \\
$\ell$ & Nonlocal length scale & m & Section 13.6 \\
\bottomrule
\end{tabularx}
\end{table}

\subsection{Open questions for Section 13}

\begin{enumerate}[label=\textbf{OQ13.\arabic*}]
\item Which applications have the strongest experimental validation across theory families, and can that validation be expressed as a single discriminating measurement rather than a set of qualitative agreements? \label{oq:13.1}
\item What is the barrier to translating generalized thermoelasticity into practical applications: the identifiability of the additional material parameters, the absence of closed-form solutions, or the lack of experimental protocols that separate the competing mechanisms? \label{oq:13.2}
\item For the nanoscale and semiconductor applications in particular, the nonlocal length scale and the memory kernel are introduced as independent parameters. Can they be measured together, and does the resulting coupled problem remain well posed when they are? \label{oq:13.3}
\item The energy-harvesting application depends on the competition between thermal and mechanical channels. Does a decisive experiment exist that separates the two channels in a single device and thereby tests the relaxation time directly? \label{oq:13.4}
\end{enumerate}

% LocalWords:
